% Оформление: ivanshemidko/algebra-notes, s2/main.tex,
% commit d52e89ab46b905aac035e556794af91aec7c86ef.
\usepackage{fontspec}
% Project-local installation keeps local and CI font files identical.
\setmainfont{cmunrm.ttf}[
  Path=fonts/,
  BoldFont=cmunbx.ttf,
  ItalicFont=cmunti.ttf,
  BoldItalicFont=cmunbi.ttf
]
\usepackage[english,russian]{babel}
\usepackage{amsmath,amssymb,mathtools,amsthm}
% Unicode maps for traditional Computer Modern math fonts (search/copy text).
\pdfvariable gentounicode=1
\protected\def\pdfglyphtounicode{\pdfextension glyphtounicode}
\input{glyphtounicode.tex}
% The standard map omits CMEX's sized delimiter glyph names.
\ExplSyntaxOn
\clist_map_inline:nn {big,Big,bigg,Bigg} {
  \pdfglyphtounicode{parenleft#1}{0028}
  \pdfglyphtounicode{parenright#1}{0029}
  \pdfglyphtounicode{braceleft#1}{007B}
  \pdfglyphtounicode{braceright#1}{007D}
  \pdfglyphtounicode{bracketleft#1}{005B}
  \pdfglyphtounicode{bracketright#1}{005D}
  \pdfglyphtounicode{floorleft#1}{230A}
  \pdfglyphtounicode{floorright#1}{230B}
  \pdfglyphtounicode{ceilingleft#1}{2308}
  \pdfglyphtounicode{ceilingright#1}{2309}
  \pdfglyphtounicode{angbracketleft#1}{27E8}
  \pdfglyphtounicode{angbracketright#1}{27E9}
  \pdfglyphtounicode{slash#1}{002F}
  \pdfglyphtounicode{backslash#1}{005C}
}
\ExplSyntaxOff
\usepackage{geometry}
\geometry{left=22mm,right=22mm,top=25mm,bottom=27mm,heightrounded}
\usepackage{enumitem}
\usepackage[most]{tcolorbox}
\tcbuselibrary{theorems,breakable,skins}
\usepackage{graphicx,caption}
\captionsetup{font=small,labelfont=bf,hypcap=false}

\definecolor{cbBlueBack}{RGB}{246,249,252}
\definecolor{cbGreenBack}{RGB}{246,251,248}
\definecolor{cbOrangeBack}{RGB}{252,249,242}
\definecolor{cbPurpleBack}{RGB}{249,247,252}
\definecolor{cbGrayBack}{RGB}{249,249,249}
\definecolor{cbBlue}{RGB}{42,87,132}
\definecolor{cbGreen}{RGB}{52,110,72}
\definecolor{cbOrange}{RGB}{140,98,0}
\definecolor{cbPurple}{RGB}{93,78,140}
\definecolor{cbGray}{RGB}{80,80,80}

\newtheoremstyle{mystatement}{8pt}{8pt}{\itshape}{}{\bfseries}{.}{\newline}{}
\newtheoremstyle{mydefinition}{8pt}{8pt}{\normalfont}{}{\bfseries}{.}{0.5em}{}
\theoremstyle{mystatement}
\newtheorem{theorem}{Теорема}[part]
\newtheorem{lemma}{Лемма}[part]
\newtheorem{cor}{Следствие}[part]
\theoremstyle{mydefinition}
\newtheorem{definition}{Определение}[part]
\newtheorem{example}{Пример}[part]
\numberwithin{equation}{part}
\counterwithin{figure}{part}
\tcbset{
  enhanced,sharp corners,boxrule=0.5pt,
  left=9pt,right=9pt,top=7pt,bottom=7pt,
  before skip=8pt,after skip=8pt,breakable
}
\tcolorboxenvironment{theorem}{colback=cbBlueBack,colframe=cbGray!55,borderline west={1.6pt}{0pt}{cbBlue}}
\tcolorboxenvironment{lemma}{colback=cbBlueBack,colframe=cbGray!55,borderline west={1.6pt}{0pt}{cbBlue}}
\tcolorboxenvironment{cor}{colback=cbPurpleBack,colframe=cbGray!55,borderline west={1.6pt}{0pt}{cbPurple}}
\tcolorboxenvironment{definition}{colback=cbGreenBack,colframe=cbGreen!55!cbGray,borderline west={1.6pt}{0pt}{cbGreen}}
\tcolorboxenvironment{example}{colback=cbOrangeBack,colframe=cbOrange!45!cbGray,borderline west={1.6pt}{0pt}{cbOrange}}

\usepackage{hyperref}
\hypersetup{
  colorlinks=true,linkcolor=cbBlue,citecolor=cbBlue,urlcolor=cbBlue,
  pdftitle={Математический анализ · III семестр},
  pdfauthor={},pdfsubject={},pdfcreator={LuaLaTeX}
}
\linespread{1.0}
\setlength{\parindent}{0pt}
\setlength{\parskip}{4pt}
\allowdisplaybreaks[1]
\providecommand{\tightlist}{\setlength{\itemsep}{0pt}\setlength{\parskip}{0pt}}

\usepackage{titlesec}
\titleformat{\part}[block]{\Large\bfseries\centering}{}{0pt}{}
\titlespacing*{\part}{0pt}{0pt}{12pt}
\counterwithin{section}{part}
\renewcommand{\thepart}{\arabic{part}}
\usepackage{tocloft}
\setlength{\cftsecnumwidth}{2.8em}
\setlength{\cftsubsecnumwidth}{3.6em}
\setcounter{tocdepth}{2}
\usepackage{fancyhdr}
\pagestyle{fancy}
\fancyhf{}
\fancyhead[L]{\nouppercase{\leftmark}}
\fancyfoot[C]{\thepage}
\let\oldpart\part
\renewcommand{\part}[1]{\oldpart{#1}\markboth{#1}{}}
\renewcommand{\sectionmark}[1]{}
\usepackage{etoolbox}
\pretocmd{\part}{\newpage}{}{}
